% experiment.tex

\chapter{实验设计与结果分析}
\echapter{Experiments and Analysis}

\section{实验环境与参数设置}
\esection{Experimental Setup}

本文实验基于修改后的 TimeMixer 项目完成。根据项目记录，主要实验在 AutoDL 云平台的 RTX 4090 GPU 环境上运行，代码基于 PyTorch 实现。当前项目中的关键结果文件保存在结果表目录中，包括完整 8 变体实验结果和各变体胜出次数统计。

实验任务为时序数据多步预测。输入序列长度统一设置为 96，预测长度分别设置为 96、192、336 和 720。模型训练默认使用代码中的 MSE 损失函数，评价时统计 MSE 和 MAE 两项指标。为保证消融实验尽量公平，在同一数据集和同一预测长度下，除对应改进模块开关外，其余主要参数保持一致。

三个改进模块通过命令行参数控制：使用 \texttt{--down\_sampling\_method conv\_ds} 启用深度可分离卷积下采样；使用 \texttt{--decomp\_method multi\_gated} 启用多核门控分解；使用 \texttt{--channel\_mixing 1} 启用通道混合。完整消融实验共包含 8 个模型变体、8 个数据集和 4 种预测长度，因此共有 $8 \times 8 \times 4 = 256$ 个模型设置。每个设置统计 MSE 和 MAE 两项指标，最终形成 64 个“数据集/预测长度/指标”评价场景。

\section{数据集介绍}
\esection{Datasets}

本文使用 8 个公开时序预测数据集，覆盖电力、气象、交通和能源等场景。数据集基本情况如表 \ref{tab:datasets} 所示。表中的变量数和采样频率来自项目记录和常用公开数据集说明，用于说明不同数据集的规模和特点。

\begin{table}[htbp]
    \centering
    \caption{实验数据集统计表}
    \label{tab:datasets}
    \begin{tabular}{cccl}
        \toprule
        数据集 & 变量数 & 采样频率 & 数据特点 \\
        \midrule
        ETTh1 / ETTh2 & 7 & 1 hour & 电力变压器油温与负载，时间粒度较粗 \\
        ETTm1 / ETTm2 & 7 & 15 min & 电力变压器数据，包含较细粒度变化 \\
        Weather & 21 & 10 min & 多气象指标，变量之间存在一定关联 \\
        ECL & 321 & 1 hour & 多用户用电量，周期性较明显 \\
        Traffic & 862 & 1 hour & 道路占用率，变量数较多且波动复杂 \\
        Solar & 137 & 10 min & 太阳能发电功率，受天气影响较大 \\
        \bottomrule
    \end{tabular}
\end{table}

从数据特点看，ETT 系列变量数较少，适合观察模型在电力变压器数据上的基本表现；Weather、ECL、Traffic 和 Solar 的变量数更多，更能体现多变量建模和模型稳定性的差异。本文不单独追求某一个数据集上的最优结果，而是从多个数据集和预测长度的整体表现分析改进模块是否稳定有效。

\section{评价指标}
\esection{Evaluation Metrics}

本文采用均方误差（Mean Squared Error, MSE）和平均绝对误差（Mean Absolute Error, MAE）作为评价指标。设真实值为 $y_i$，预测值为 $\hat{y}_i$，样本总数为 $n$，则 MSE 定义为：
\begin{equation}
    \mathrm{MSE} = \frac{1}{n}\sum_{i=1}^{n}(y_i-\hat{y}_i)^2.
\end{equation}
MAE 定义为：
\begin{equation}
    \mathrm{MAE} = \frac{1}{n}\sum_{i=1}^{n}|y_i-\hat{y}_i|.
\end{equation}

MSE 对较大误差更敏感，能够反映模型是否出现明显预测偏差；MAE 直接表示平均绝对误差，更容易从整体误差大小进行理解。两项指标均为数值越低表示预测效果越好。本文在分析实验结果时同时观察 MSE 和 MAE，避免只依据单一指标作出结论。

\section{主实验结果}
\esection{Main Results}

表 \ref{tab:main_results} 给出了 8 个模型变体在 8 个数据集和 4 种预测长度上的完整结果。由于结果表较宽，表中保留了 MSE/MAE 并列展示的形式，粗体表示同一数据集和同一预测长度下对应指标的最优值。

\begin{table}[htbp]
\centering
\caption{时序数据预测实验结果（MSE/MAE，最优结果加粗）}
\label{tab:main_results}
\resizebox{\textwidth}{!}{%
\begin{tabular}{cc|cc|cc|cc|cc|cc|cc|cc|cc}
\toprule
\multicolumn{2}{c|}{} & \multicolumn{2}{c|}{M0} & \multicolumn{2}{c|}{M1} & \multicolumn{2}{c|}{M2} & \multicolumn{2}{c|}{M3} & \multicolumn{2}{c|}{M12} & \multicolumn{2}{c|}{M13} & \multicolumn{2}{c|}{M23} & \multicolumn{2}{c|}{M4} \\
数据集 & 预测长度 & MSE & MAE & MSE & MAE & MSE & MAE & MSE & MAE & MSE & MAE & MSE & MAE & MSE & MAE & MSE & MAE \\
\midrule
\multirow{4}{*}{ETTh1} & 96 & 0.3865 & 0.4029 & 0.3929 & 0.4086 & \textbf{0.3809} & \textbf{0.3989} & 0.4129 & 0.4180 & 0.3894 & 0.4036 & 0.3914 & 0.4147 & 0.3843 & 0.4039 & 0.4138 & 0.4191 \\
 & 192 & 0.4420 & \textbf{0.4294} & 0.4458 & 0.4349 & \textbf{0.4395} & 0.4350 & 0.4602 & 0.4476 & 0.4461 & 0.4343 & 0.4484 & 0.4388 & 0.4401 & 0.4311 & 0.4648 & 0.4462 \\
 & 336 & 0.5135 & 0.4704 & 0.6166 & 0.5426 & \textbf{0.4725} & \textbf{0.4451} & 0.5117 & 0.4660 & 0.4886 & 0.4615 & 0.5863 & 0.5199 & 0.5083 & 0.4711 & 0.4801 & 0.4538 \\
 & 720 & \textbf{0.5075} & \textbf{0.4806} & 0.5309 & 0.5042 & 0.5095 & 0.4905 & 0.5114 & 0.4854 & 0.5388 & 0.5037 & 0.5560 & 0.5119 & 0.5803 & 0.5213 & 0.5154 & 0.4936 \\
\midrule
\multirow{4}{*}{ETTh2} & 96 & \textbf{0.2908} & 0.3416 & 0.2934 & 0.3443 & 0.2957 & \textbf{0.3415} & 0.3219 & 0.3622 & 0.2995 & 0.3499 & 0.3117 & 0.3568 & 0.3194 & 0.3588 & 0.3174 & 0.3636 \\
 & 192 & 0.3772 & \textbf{0.3947} & 0.3929 & 0.4067 & \textbf{0.3761} & 0.3982 & 0.3974 & 0.4071 & 0.4036 & 0.4106 & 0.3764 & 0.3969 & 0.3984 & 0.4099 & 0.4261 & 0.4216 \\
 & 336 & \textbf{0.4173} & \textbf{0.4314} & 0.4299 & 0.4399 & 0.4234 & 0.4388 & 0.4385 & 0.4431 & 0.4411 & 0.4428 & 0.5793 & 0.5371 & 0.4305 & 0.4358 & 0.4306 & 0.4402 \\
 & 720 & 0.4638 & 0.4635 & 0.4956 & 0.4781 & 0.4696 & 0.4702 & \textbf{0.4367} & \textbf{0.4537} & 0.6511 & 0.5736 & 0.9467 & 0.6884 & 0.5019 & 0.4891 & 0.4847 & 0.4786 \\
\midrule
\multirow{4}{*}{ETTm1} & 96 & 0.3317 & 0.3669 & 0.3171 & 0.3572 & 0.3188 & 0.3577 & 0.3370 & 0.3737 & 0.3266 & 0.3586 & \textbf{0.3099} & \textbf{0.3530} & 0.3557 & 0.3849 & 0.3402 & 0.3726 \\
 & 192 & 0.3627 & 0.3821 & 0.3869 & 0.4018 & \textbf{0.3577} & 0.3812 & 0.3829 & 0.4019 & 0.3614 & \textbf{0.3809} & 0.3842 & 0.3995 & 0.3878 & 0.4025 & 0.3829 & 0.3995 \\
 & 336 & \textbf{0.3889} & \textbf{0.4027} & 0.3973 & 0.4074 & 0.4004 & 0.4075 & 0.4058 & 0.4149 & 0.4246 & 0.4266 & 0.4239 & 0.4216 & 0.4010 & 0.4128 & 0.3968 & 0.4081 \\
 & 720 & \textbf{0.4616} & \textbf{0.4427} & 0.4754 & 0.4599 & 0.4691 & 0.4510 & 0.4738 & 0.4537 & 0.4651 & 0.4445 & 0.5126 & 0.4821 & 0.4683 & 0.4537 & 0.4865 & 0.4622 \\
\midrule
\multirow{4}{*}{ETTm2} & 96 & 0.1750 & \textbf{0.2567} & 0.1751 & 0.2596 & 0.1748 & 0.2570 & 0.1822 & 0.2637 & 0.1903 & 0.2726 & 0.1907 & 0.2765 & \textbf{0.1743} & 0.2586 & 0.1795 & 0.2652 \\
 & 192 & \textbf{0.2389} & \textbf{0.3005} & 0.2587 & 0.3180 & 0.2403 & 0.3012 & 0.2496 & 0.3077 & 0.2512 & 0.3095 & 0.2399 & 0.3060 & 0.2435 & 0.3030 & 0.2582 & 0.3198 \\
 & 336 & \textbf{0.3040} & 0.3416 & 0.4039 & 0.4019 & 0.3050 & \textbf{0.3410} & 0.3117 & 0.3505 & 0.3324 & 0.3645 & 0.3505 & 0.3736 & \textbf{0.3040} & 0.3439 & 0.3043 & 0.3482 \\
 & 720 & 0.3950 & 0.4000 & 0.6187 & 0.5150 & \textbf{0.3936} & \textbf{0.3963} & 0.4070 & 0.4006 & 0.6591 & 0.5502 & 0.4222 & 0.4111 & 0.4182 & 0.4117 & 0.4319 & 0.4195 \\
\midrule
\multirow{4}{*}{Weather} & 96 & 0.1624 & 0.2100 & 0.1643 & 0.2113 & 0.1661 & 0.2138 & \textbf{0.1562} & \textbf{0.2076} & 0.1702 & 0.2155 & 0.1608 & 0.2135 & 0.1596 & 0.2086 & 0.1630 & 0.2141 \\
 & 192 & \textbf{0.2086} & 0.2527 & 0.2115 & \textbf{0.2522} & 0.2102 & 0.2529 & 0.2145 & 0.2590 & 0.2116 & 0.2569 & \textbf{0.2086} & 0.2569 & 0.2091 & 0.2555 & 0.2170 & 0.2652 \\
 & 336 & 0.2637 & 0.2933 & 0.2698 & 0.2970 & 0.2650 & \textbf{0.2914} & 0.2677 & 0.2979 & 0.2630 & 0.2938 & 0.2760 & 0.3086 & 0.2672 & 0.2954 & \textbf{0.2600} & 0.2935 \\
 & 720 & 0.3454 & 0.3475 & 0.3360 & 0.3455 & 0.3455 & 0.3471 & 0.3395 & 0.3448 & 0.3286 & 0.3420 & 0.3322 & \textbf{0.3397} & 0.3447 & 0.3462 & \textbf{0.3279} & 0.3456 \\
\midrule
\multirow{4}{*}{ECL} & 96 & 0.1552 & 0.2469 & 0.1555 & 0.2485 & 0.1545 & 0.2459 & 0.1569 & 0.2579 & \textbf{0.1475} & \textbf{0.2435} & 0.1590 & 0.2585 & 0.1564 & 0.2527 & 0.1641 & 0.2633 \\
 & 192 & 0.1677 & 0.2585 & \textbf{0.1645} & 0.2572 & 0.1672 & \textbf{0.2571} & 0.1741 & 0.2739 & 0.1672 & 0.2594 & 0.1969 & 0.2924 & 0.1729 & 0.2707 & 0.1759 & 0.2733 \\
 & 336 & 0.1846 & 0.2755 & 0.1906 & 0.2826 & \textbf{0.1838} & \textbf{0.2739} & 0.1898 & 0.2919 & 0.1991 & 0.2875 & 0.1922 & 0.2912 & 0.1842 & 0.2794 & 0.1879 & 0.2884 \\
 & 720 & \textbf{0.2241} & 0.3101 & 0.2379 & 0.3248 & 0.2254 & \textbf{0.3098} & 0.2334 & 0.3248 & 0.2551 & 0.3367 & 0.2488 & 0.3346 & 0.2273 & 0.3271 & 0.2285 & 0.3199 \\
\midrule
\multirow{4}{*}{Traffic} & 96 & 0.5181 & 0.3505 & 0.5202 & 0.3469 & \textbf{0.4621} & \textbf{0.2844} & 0.5530 & 0.3783 & 0.5387 & 0.3570 & 0.4993 & 0.3418 & 0.4846 & 0.3331 & 0.5296 & 0.3634 \\
 & 192 & \textbf{0.4655} & \textbf{0.2914} & 0.5361 & 0.3214 & 0.4682 & \textbf{0.2914} & 0.5388 & 0.3583 & 0.4668 & 0.3191 & 0.5674 & 0.4153 & 0.5151 & 0.3641 & 0.5005 & 0.3473 \\
 & 336 & 0.5013 & 0.3022 & 0.5288 & 0.3423 & \textbf{0.4837} & \textbf{0.2979} & 0.5524 & 0.3477 & 0.5587 & 0.3454 & 0.5513 & 0.3544 & 0.5053 & 0.3382 & 0.5385 & 0.3635 \\
 & 720 & 0.5528 & 0.3132 & 0.5865 & 0.3625 & \textbf{0.5085} & \textbf{0.3112} & 0.5626 & 0.3754 & 0.5370 & 0.3563 & 0.5689 & 0.3697 & 0.6103 & 0.4015 & 0.5431 & 0.3449 \\
\midrule
\multirow{4}{*}{Solar} & 96 & 0.2014 & 0.2658 & 0.1902 & 0.2501 & 0.2136 & 0.2914 & 0.1922 & 0.2846 & \textbf{0.1871} & \textbf{0.2344} & 0.2063 & 0.2681 & 0.1893 & 0.2508 & 0.1961 & 0.2934 \\
 & 192 & 0.2435 & 0.3119 & 0.2217 & \textbf{0.2516} & 0.2632 & 0.3393 & 0.2441 & 0.2667 & 0.2212 & 0.2693 & 0.2197 & 0.3078 & 0.2196 & 0.2996 & \textbf{0.1993} & 0.2734 \\
 & 336 & 0.2429 & 0.2930 & 0.2365 & 0.2893 & 0.2533 & 0.3317 & 0.2355 & 0.3078 & 0.2248 & \textbf{0.2617} & 0.2248 & 0.2875 & 0.2307 & 0.2862 & \textbf{0.2222} & 0.2939 \\
 & 720 & 0.2448 & 0.2914 & 0.2287 & 0.2948 & 0.2694 & 0.3221 & 0.2411 & 0.3157 & 0.2235 & 0.2886 & 0.2416 & 0.3194 & 0.2175 & \textbf{0.2757} & \textbf{0.2159} & 0.2841 \\
\bottomrule
\end{tabular}}
\end{table}


从完整结果可以看出，原始 TimeMixer 基线 M0 仍然具有较强竞争力，在不少数据集和预测长度上可以取得最优或接近最优的结果。这说明原始模型的多尺度分解混合结构本身较为有效，也说明本文改进不应简单表述为“所有模块都显著提升”。相比之下，M2 在多个场景中表现稳定，尤其在 Traffic 等周期结构复杂、变量数较多的数据集上取得了较明显优势。M1 和 M3 在部分数据集上有局部收益，但结果波动较大。M4 虽然包含全部三个改进模块，但并没有在整体上全面超过其他变体。

\section{单模块消融实验分析}
\esection{Single-module Ablation Analysis}

为观察单个模块的独立作用，可以重点比较 M0、M1、M2 和 M3。M1 只替换下采样模块，M2 只替换分解模块，M3 只加入通道混合模块。

从结果看，M2 是三种单模块改进中最稳定的一项。它在 ETTh1、ETTm1、ETTm2、ECL 和 Traffic 等数据集的若干预测长度上取得最优结果。尤其在 Traffic 数据集上，M2 在多个预测长度下均取得较好 MSE 或 MAE，说明多核门控分解对复杂周期结构具有较好的适应能力。造成这一现象的原因可能是 Traffic 数据中既有日常出行周期，也有局部突发拥堵，单一窗口分解难以同时兼顾，而多核窗口可以提供更灵活的趋势估计。

M1 的表现相对不稳定。深度可分离卷积下采样可以让模型学习下采样权重，有机会保留局部突变信息，但它也可能引入额外噪声或改变原始序列的平滑结构。对于 ETTm 等细粒度数据，M1 在部分预测长度上能取得较好结果；但在一些较平滑或预测步长较大的场景中，它未必优于平均池化。这说明可学习下采样并不是无条件有效，需要结合数据特征进一步调节。

M3 的通道混合模块也表现出一定局部性。该模块在 Weather 等变量相关性较强的数据上有时能取得较好表现，但整体稳定性不足。一个可能原因是当前通道混合只发生在输入阶段，且只进行一次浅层低秩映射，对复杂变量关系的表达能力有限。同时，在通道独立主干结构中加入额外变量交互，也可能改变输入分布，使后续分解和混合过程更难优化。

\section{组合消融实验分析}
\esection{Combined Ablation Analysis}

除了单模块实验，本文还补充了 M12、M13 和 M23 三个两两组合变体，并与 M4 完整模型进行比较。这样可以进一步观察模块之间是否存在互补关系或耦合干扰。

表 \ref{tab:wins} 给出了各模型变体在 64 个评价场景中的胜出次数。M2 的总胜出次数为 22 次，是所有变体中最高的；M0 获得 19 次，说明原始模型仍然是较强基线；M4 仅获得 5 次，低于 M2 和 M0，也低于部分组合变体的局部表现。

\begin{table}[htbp]
    \centering
    \caption{各模型变体最优次数统计}
    \label{tab:wins}
    \begin{tabular}{cccc}
        \toprule
        变体 & MSE最优次数 & MAE最优次数 & 总计 \\
        \midrule
        M0 & 10 & 9 & 19 \\
        M1 & 1 & 2 & 3 \\
        M2 & 10 & 12 & 22 \\
        M3 & 2 & 2 & 4 \\
        M12 & 2 & 4 & 6 \\
        M13 & 1 & 2 & 3 \\
        M23 & 1 & 1 & 2 \\
        M4 & 5 & 0 & 5 \\
        \bottomrule
    \end{tabular}
\end{table}

从组合结果看，多个模块同时加入并不一定带来线性叠加收益。M12、M13 和 M23 在部分场景下优于 M4，说明完整模型中不同模块之间可能存在耦合干扰。项目记录中还统计了相对 M4 的去模块观察：M12 在 MSE 上有 14/32 个场景优于 M4，在 MAE 上有 19/32 个场景优于 M4；M13 分别为 11/32 和 15/32；M23 分别为 18/32 和 20/32。这说明去掉某些模块后，模型反而可能在部分场景中更稳定。

造成这一现象的原因可能有三点。首先，M1 的可学习下采样和 M2 的多核分解都在处理输入特征与趋势信息，二者功能上可能存在一定重叠。其次，M3 改变了变量维度的信息分布，可能使后续分解模块面对的输入特征与原始设定不同。最后，三个模块同时加入后参数和优化路径都变得更复杂，在固定训练设置下未必能够充分发挥作用。因此，本文不将 M4 写作“最终最优模型”，而是将其视为结构最完整的集成变体，并如实分析其不足。

\section{实验结果讨论}
\esection{Discussion}

综合实验结果，本文可以得到以下几点认识。

第一，多核门控分解是本文最主要、最稳定的有效改进。它直接作用于 TimeMixer 的核心分解机制，使模型能够利用多个窗口估计趋势项，并通过门控权重进行自适应融合。相比下采样和通道混合，分解模块更接近 TimeMixer 的结构瓶颈，因此对整体性能影响更明显。

第二，原始 TimeMixer 仍然具有较强竞争力。M0 获得 19 次最优，说明原始多尺度分解混合结构本身已经比较有效。本文的改进结果不能简单理解为“新模块一定提升”，而应理解为对原模型不同环节的探索和验证。

第三，深度可分离下采样和通道混合具有一定探索价值，但稳定性不足。M1 能够提供可学习下采样能力，M3 能够引入变量间交互，但二者在当前实现中没有形成稳定优势。后续如果继续改进，可以考虑为 M1 增加更明确的平滑约束，或将 M3 从单次输入混合扩展为分层通道交互。

第四，模块融合方式仍需进一步研究。M4 结构最完整，但实验显示简单叠加并非总能带来更优结果。这说明在轻量级时序预测模型中，不同结构模块之间的耦合关系需要谨慎处理。对于本科毕业设计而言，这一结论也比较符合实际：模型改进不只是添加模块，还需要通过消融实验判断哪些改进真正有效。
